\prefacesection{Agradecimentos}
Gostaria de começar estes agradecimentos relembrando o quanto cada pessoa aqui citada foi essencial na construção deste trabalho e na minha jornada acadêmica.

Primeiramente, meu muito obrigado ao Professor Edson Borin, meu orientador.
Ele foi a primeira pessoa que realmente acreditou e confiou na minha capacidade acadêmica, mesmo eu vindo de uma área totalmente diferente.
Foi graças à confiança dele que me senti seguro para mergulhar no mundo da computação, da academia e da ciência.
Ele não só me apresentou novos horizontes de pesquisa, como também me apoiou de forma incondicional em cada passo.
Sua paciência, sua experiência e sua visão foram cruciais para que eu pudesse amadurecer meus conhecimentos e avançar com segurança no desenvolvimento desta tese.

Em segundo lugar, gostaria de expressar minha profunda gratidão ao Professor Rodolfo Azevedo, que abriu as portas para mim em um momento em que muitos outros caminhos se fechavam.
Ele enxergou potencial e me ajudou a vislumbrar um futuro, mesmo quando eu estava bastante perdido e sem clareza de como poderia moldar minha carreira.
Foi nas mentorias com ele que descobri modos de organizar minhas ideias, articular objetivos e estruturar um plano que fizesse sentido para mim tanto profissionalmente quanto academicamente.
Ele foi fundamental para que eu pudesse me preparar e ingressar no mestrado de computação.
Posso afirmar que sua clareza, empatia e generosidade foram essenciais para que eu pudesse dar os primeiros passos em direção ao futuro que tanto almejava.

Também gostaria de agradecer ao professor Carlos César de Araújo.
Tive o privilégio de contar com sua orientação durante meu preparo para ingressar no mestrado.
Sua paciência e dedicação foram fundamentais para me ensinar conceitos fundamentais de lógica e matemática que me faltavam no início da minha jornada.
Ele me ajudou a desenvolver uma base sólida, o que foi crucial para que eu pudesse avançar em minha formação acadêmica.

Também não posso deixar de agradecer ao Samuel Goto, que teve uma paciência enorme comigo e me acompanhou por um longo período.
Ele me guiou em mentorias valiosas, compartilhando não só seus conhecimentos técnicos, mas também oferecendo conselhos pessoais e de carreira que me ajudaram a crescer de forma integral.
A maneira como ele dedicou tempo para ouvir minhas inseguranças e me mostrar novos caminhos fez toda a diferença nesta trajetória.

Por fim, mas não menos importante, estendo meus agradecimentos à Petrobras, pela oportunidade e pelo espaço que me foi concedido para desenvolver esta pesquisa.
Sem o suporte e a estrutura oferecida pela equipe, este trabalho não teria saído do papel.
O apoio do time e a confiança no projeto foram fundamentais para que eu pudesse mergulhar de cabeça nos estudos e levar esta tese até o final.

A cada um de vocês, meu mais sincero obrigado.